% Lösungen Einheit 2 von 3 – Zehnerpotenzen (zusammenbau v0.1)
\einheitenkopf{Einheit 2 von 3 · Zehnerpotenzen}

\erg{19}{a) größer als eins – die Hochzahl ist positiv, das Komma wandert nach rechts \quad b) $10\,000$ \quad c) Milliarde – neun Nullen \quad d) $5\,300$ \quad e) $2\,900\,000$ \quad f) $4\,800$ \quad g) $630\,000$ – das Komma wandert fünf Stellen nach rechts \quad h) $5{,}8 \cdot 10^{4}$ (Hochzahl $= 4$) \quad i) $4$ – das Komma wandert vier Stellen \quad j) $86\,000\,000\,000$ \quad k) $0{,}0036$ \quad l) $5{,}2 \cdot 10^{-3}$ (Hochzahl $= -3$) \quad m) $93\,000 < 410\,000 < 2\,800\,000$, also $9{,}3 \cdot 10^{4} < 4{,}1 \cdot 10^{5} < 2{,}8 \cdot 10^{6}$ \quad n) $6{,}2 \cdot 10^{7} = 62\,000\,000$ \quad o) $149\,597\,871$ km $\approx 1{,}50 \cdot 10^{8}$ km \quad p) $2 \cdot 10^{-3}$ mm $= 0{,}002$ mm; $1 : 0{,}002 = 500$ Bakterien \quad q) $10^{4+3} = 10^{7}$ (Hochzahl $= 7$) \quad r) $5$ – das Komma wandert fünf Stellen}
\erg{20}{a) Nullen gezählt statt Stellen: Das Komma wandert fünf Stellen, also $920\,000 = 9{,}2 \cdot 10^{5}$ (Hochzahl $= 5$) \quad b) Beide sind $260\,000$: Ist der Faktor zehnmal so groß, muss die Hochzahl um eins kleiner sein.}
